\subsection{CONDITIONS FOR AN INVERSE LIMIT TO BE NON-EMPTY}

In this subsection we shall give the two most frequently used sufficient conditions for an inverse limit to be non-empty (see also Exercise 5).

PROPOSITION 5. Let \((\mathbf{E}_{\alpha}, f_{\alpha \beta})\) be an inverse system of sets relative to a directed set I which has a countable cofinal subset, and suppose furthermore that the \(f_{\alpha \beta}\) are surjective. Then, if \(\mathbf{E} = \lim \mathbf{E}_{\alpha}\) , the canonical mapping \(f_{x}: \mathbf{E} \to \mathbf{E}_{\alpha}\) is surjective for each \(\alpha \in \mathbf{I}\) (and, a fortiori, \(\mathbf{E}\) is not empty provided that none of the \(\mathbf{E}_{\alpha}\) is empty).

Let \((\alpha_{n})\) be a sequence of elements of I which form a cofinal subset of I. Since I is directed, we can define a sequence \((\beta_{n})\) of elements of I inductively by the conditions \(\beta_{0} = \alpha_{0}, \beta_{n} \geqslant \beta_{i}\) for \(i < n\) and \(\beta_{n} \geqslant \alpha_{n}\) . Clearly the sequence \((\beta_{n})\) is increasing and forms a cofinal subset of I. In view of Proposition 1 of no. 1 and the relations \(f_{\alpha} = f_{\alpha \beta_{n}} \circ f_{\beta_{n}}\) for \(\alpha \leqslant \beta_{n}\) , we need only prove the Proposition for the case I = N. Moreover, it is clear that it suffices to prove that \(f_{0}\) is surjective. Let \(x_{0} \in \mathrm{E}_{0}\) . Define \(x_{n} \in \mathrm{E}_{n}\) (for \(n \geqslant 1\) ) inductively to be an element of the set \(\overline{f}_{n-1,n}(x_{n-1})\) , which is possible because the latter set is non-empty by hypothesis. We then show by induction on \(n - m\) that \(x_{m} = f_{mn}(x_{n})\) for \(m \leqslant n\) , and it follows that \(x = (x_{n})\) belongs to E.

The second criterion concerns inverse systems \((\mathbf{E}_{\alpha}, f_{\alpha\beta})\) relative to an index set I such that for each \(\alpha \in I\) we are given a set \(S_{\alpha}\) of subsets of \(E_{\alpha}\) which satisfy the following conditions:

\begin{enumerate}
\def\labelenumi{(\roman{enumi})}
\tightlist
\item
  Every intersection of sets belonging to \(\mathfrak{S}_{\alpha}\) also belongs to \(\mathfrak{S}_{\alpha}\) .
\end{enumerate}

It follows in particular (by considering the intersection of the empty family) that \(\mathbf{E}_{\alpha} \in \mathfrak{S}_{\alpha}\) .

\begin{enumerate}
\def\labelenumi{(\roman{enumi})}
\setcounter{enumi}{1}
\tightlist
\item
  If a set of subsets \(\mathfrak{F} \subset \mathfrak{S}_{\alpha}\) is such that every finite intersection of sets belonging to \(\mathfrak{F}\) is non-empty, then \(\bigcap_{\mathbf{M} \in \mathfrak{F}} \mathbf{M}\) is non-empty.
\end{enumerate}

In view of (i), it is clear that (ii) is equivalent to the following condition :

(ii)' If \(\mathfrak{G} \subset \mathfrak{S}_{\alpha}\) is left directed (with respect to inclusion) and does not contain the empty set, then \(\bigcap_{M \in \mathfrak{G}} M\) is non-empty.

THEOREM 1. Suppose that I is directed, that the sets \(\mathfrak{S}_{\alpha}\) satisfy conditions (i) and (ii), and that the inverse system \((\mathrm{E}_{\alpha}, f_{\alpha\beta})\) has the following properties:

\begin{enumerate}
\def\labelenumi{(\roman{enumi})}
\setcounter{enumi}{2}
\item
  For each pair of indices \(\alpha, \beta\) such that \(\alpha \leqslant \beta\) , and each \(x_{\alpha} \in \mathbf{E}_{\alpha}\) , we have \(\overline{f}_{\alpha\beta}(x_{\alpha}) \in \mathfrak{S}_{\beta}\) .
\item
  For each pair of indices \(\alpha, \beta\) such that \(\alpha \leqslant \beta\) , and each \(\mathbf{M}_{\beta} \in \mathfrak{S}_{\beta}\) , we have \(f_{\alpha \beta}(\mathbf{M}_{\beta}) \in \mathfrak{S}_{\alpha}\) .
\end{enumerate}

Let \(\mathbf{E} = \varprojlim \mathbf{E}_{\alpha}\) and let \(f_{\alpha}: \mathbf{E} \to \mathbf{E}_{\alpha}\) be the canonical mapping for each \(\alpha \in \mathbf{I}\) . Then:

\begin{enumerate}
\def\labelenumi{(\alph{enumi})}
\tightlist
\item
  For each \(\alpha \in I\) we have
\end{enumerate}

\[
f _ {\alpha} (\mathbf {E}) = \bigcap_ {\beta \geqslant \alpha} f _ {\alpha \beta} (\mathbf {E} _ {\beta}).\tag{19}
\]

\begin{enumerate}
\def\labelenumi{(\alph{enumi})}
\setcounter{enumi}{1}
\tightlist
\item
  If \(\mathbf{E}_{\alpha}\) is non-empty for each \(\alpha \in \mathbf{I}\) , then \(\mathbf{E}\) is non-empty.
\end{enumerate}

Let \(\Sigma\) be the set of families \(\mathfrak{A} = (\mathrm{A}_{\alpha})_{\alpha \in \mathbf{I}}\) which satisfy the following conditions:

\begin{enumerate}
\def\labelenumi{(\arabic{enumi})}
\setcounter{enumi}{19}
\tightlist
\item
\end{enumerate}

\[
\mathrm{A} _ {\alpha} \neq \emptyset \quad \text{and} \quad \mathrm{A} _ {\alpha} \in \mathfrak {S} _ {\alpha} \quad \text{for all} \quad \alpha \in \mathrm{I};\tag{21}
\]

\[
f _ {\alpha \beta} (\mathbf {A} _ {\beta}) \subset \mathbf {A} _ {\alpha} \quad \text{whenever} \quad \alpha \leqslant \beta .
\]

If \(\mathfrak{A} = (\mathrm{A}_{\alpha})\) and \(\mathfrak{A}' = (\mathrm{A}_{\alpha}')\) are any two elements of \(\Sigma\) , let the relation \(\mathfrak{A} \leqslant \mathfrak{A}'\) mean that \(\mathrm{A}_{\alpha} \supset \mathrm{A}_{\alpha}'\) for all \(\alpha\) . Clearly \(\Sigma\) is ordered by this relation.

\begin{enumerate}
\def\labelenumi{(\arabic{enumi})}
\item
  Let us first show that the ordered set \(\Sigma\) is inductive. Let \(L\) be a totally ordered set and let \(\lambda \to \mathfrak{A}^{\lambda} = (A_{\alpha}^{\lambda})_{\alpha \in I}\) be a strictly increasing mapping of \(L\) into \(\Sigma\) . For each \(\alpha \in I\) , put \(B_{\alpha} = \bigcap_{\lambda \in L} A_{\alpha}^{\lambda}\) . Then it is immediately seen that the family \(\mathfrak{B} = (B_{\alpha})_{\alpha \in I}\) satisfies (21); by reason of (i) and (ii), it also satisfies (20), hence belongs to \(\Sigma\) ; and it is clear that \(\mathfrak{B}\) is an upper bound of the set of the \(\mathfrak{A}^{\lambda}\) .
\item
  Let \(\mathfrak{A} = (\mathrm{A}_{\alpha})\) be a maximal element of \(\Sigma\) . We shall show that \(\mathrm{A}_{\alpha} = f_{\alpha \beta}(\mathrm{A}_{\beta})\) whenever \(\alpha \leqslant \beta\) . Let \(\mathrm{A}_{\alpha}' = \bigcap_{\beta \geqslant \alpha} f_{\alpha \beta}(\mathrm{A}_{\beta})\) for all \(\alpha \in \mathbf{I}\) , and let us show that \(\mathfrak{A}' = (\mathrm{A}_{\alpha}')\) belongs to \(\Sigma\) . Note first that if \(\alpha \leqslant \beta \leqslant \gamma\) , we have \(f_{\alpha \gamma}(\mathrm{A}_{\gamma}) = f_{\alpha \beta}(f_{\beta \gamma}(\mathrm{A}_{\gamma})) \subset f_{\alpha \beta}(\mathrm{A}_{\beta})\) by (21); moreover, \(f_{\alpha \beta}(\mathrm{A}_{\beta}) \in \mathfrak{S}_{\alpha}\) by (iv), and \(f_{\alpha \beta}(\mathrm{A}_{\beta}) \neq \emptyset\) by (20). Hence conditions (i) and (ii) show that \(\mathfrak{A}'\) satisfies (20). Also, if \(\alpha \leqslant \beta\) , we have
\end{enumerate}

\[
f _ {\alpha \beta} (\mathbf {A} _ {3} ^ {\prime}) \subset \bigcap_ {\gamma \geqslant \beta} f _ {\alpha \beta} (f _ {\beta \gamma} (\mathbf {A} _ {\gamma})) = \bigcap_ {\gamma \geqslant \beta} f _ {\alpha \gamma} (\mathbf {A} _ {\gamma});
\]

and for each \(\delta \geqslant \alpha\) there exists \(\gamma \in I\) such that \(\gamma \geqslant \delta\) and \(\gamma \geqslant \beta\) , so that \(f_{\alpha \gamma}(A_{\gamma}) \subset f_{\alpha \delta}(A_{\delta})\) and consequently

\[
\bigcap_ {\gamma \geqslant \beta} f _ {\alpha \gamma} (\mathrm{A} _ {\gamma}) = \bigcap_ {\delta \geqslant \alpha} f _ {\alpha \delta} (\mathrm{A} _ {\delta}) = \mathrm{A} _ {\alpha} ^ {\prime}.
\]

Hence \(\mathfrak{A}'\) satisfies (21) and therefore belongs to \(\Sigma\) . Since \(\mathbf{A}_{\alpha}' \subset \mathbf{A}_{\alpha}\) for all \(\alpha\) , the maximality of \(\mathfrak{A}\) in \(\Sigma\) implies \(\mathfrak{A}' = \mathfrak{A}\) , and our assertion is proved.

\begin{enumerate}
\def\labelenumi{(\arabic{enumi})}
\setcounter{enumi}{2}
\tightlist
\item
  We shall establish next that if \(\mathfrak{A} = (\mathbf{A}_{\alpha})\) is a maximal element of \(\Sigma\) , then each of the \(\mathbf{A}_{\alpha}\) consists of a single element. Let \(x_{\alpha} \in \mathbf{A}_{\alpha}\) . For each \(\beta \geqslant \alpha\) , put \(\mathbf{B}_{\beta} = \mathbf{A}_{\beta} \cap f_{\alpha \beta}(x_{\alpha})\) ; if \(\beta\) is not \(\geqslant \alpha\) , put \(\mathbf{B}_{\beta} = \mathbf{A}_{\beta}\) ; then we shall see that \(\mathfrak{B} = (\mathbf{B}_{\beta})\) belongs to \(\Sigma\) . If \(\beta\) is not \(\geqslant \alpha\) , the relation \(\beta \leqslant \gamma\) implies \(f_{\beta \gamma}(\mathbf{B}_{\gamma}) \subset f_{\beta \gamma}(\mathbf{A}_{\gamma}) \subset \mathbf{A}_{\beta} = \mathbf{B}_{\beta}\) . If, on the other hand, \(\alpha \leqslant \beta \leqslant \gamma\) , then since
\end{enumerate}

\[
\overline {{f}} _ {\alpha \gamma} ^ {- 1} (x _ {\alpha}) = \overline {{f}} _ {\beta \gamma} ^ {- 1} (\overline {{f}} _ {\alpha \beta} ^ {- 1} (x _ {\alpha})),
\]

we have

\[
f _ {\beta \gamma} (\overline {{f _ {\alpha \gamma} ^ {- 1}}} (x _ {\alpha})) \subset \overline {{f _ {\alpha \beta} ^ {- 1}}} (x _ {\alpha}),
\]

and since \(f_{\beta \gamma}(\mathbf{A}_{\gamma}) \subset \mathbf{A}_{\beta}\) , we again have \(f_{\beta \gamma}(\mathbf{B}_{\gamma}) \subset \mathbf{B}_{\beta}\) , so that the family \(\mathfrak{B}\) satisfies (21). Since \(\mathbf{A}_{\alpha} = f_{\alpha \beta}(\mathbf{A}_{\beta})\) whenever \(\alpha \leqslant \beta\) by part (2) of the proof, it is clear that \(B_{\beta} \neq \varnothing\) for all \(\beta \in I\) . Finally, by virtue of (i) and (iii), we have \(B_{\beta} \in S_{\beta}\) for all \(\beta \in I\) , and hence \(B \in \Sigma\) . Since \(B_{\beta} \subset A_{\beta}\) for all \(\beta \in I\) , the maximality of A implies that \(B_{\beta} = A_{\beta}\) for all \(\beta\) , and in particular \(A_{\alpha} = \{x_{\alpha}\}\) .

\begin{enumerate}
\def\labelenumi{(\arabic{enumi})}
\setcounter{enumi}{3}
\tightlist
\item
  We are now in a position to prove Theorem 1. Let us start with (a). We have
\end{enumerate}

\[
f _ {\alpha} (\mathbf {E}) \subset \bigcap_ {\beta \geqslant \alpha} f _ {\alpha \beta} (\mathbf {E} _ {\beta}).
\]

Conversely, let

\[
x _ {\alpha} \in \bigcap_ {\beta \geqslant x} f _ {\alpha \beta} (\mathrm{E} _ {\beta}),
\]

and put

\[
\mathbf {B} _ {\beta} = \overline {{f}} _ {\alpha \beta} ^ {1} (x _ {\alpha})
\]

if \(\beta \geqslant \alpha\) , and \(B_{\beta} = E_{\beta}\) otherwise. By the definition of \(x_{\alpha}\) , the \(B_{\beta}\) are non-empty, and we have \(B_{\beta} \in \mathfrak{S}_{\beta}\) for all \(\beta \in I\) by virtue of (iii) and (i); moreover, it is evident that \(f_{\beta \gamma}(B_{\gamma}) \subset B_{\beta}\) whenever \(\beta \leqslant \gamma\) . Hence \(\mathfrak{B} = (B_{\beta}) \in \Sigma\) . Let \(\mathfrak{A} = (A_{\alpha})\) be a maximal element of \(\Sigma\) such that \(\mathfrak{A} \geqslant \mathfrak{B}\) (the existence of \(\mathfrak{A}\) follows from (1) and §2, no. 4, Theorem 2, Corollary 1). Since, by (3), \(A_{\beta}\) is of the form \(\{y_{\beta}\}\) for all \(\beta \in I\) , it follows that \(y = (y_{\beta})\) belongs to \(E\) , and \(f_{\alpha}(y) = y_{\alpha} = x_{\alpha}\) by definition.

Finally, (a) implies (b). We may as well assume that I is not empty (otherwise there is nothing to prove). The hypothesis that the \(E_{\alpha}\) are non-empty implies that \(f_{\alpha\beta}(E_{\beta}) \neq \emptyset\) for all \(\beta \geqslant \alpha\) . Since the \(f_{\alpha\beta}(E_{\beta})\) , for fixed \(\alpha\) and variable \(\beta \geqslant \alpha\) , form a left directed set of subsets of \(E_{\alpha}\) belonging to \(S_{\alpha}\) , condition (ii)' proves that

\[
\bigcap_ {\beta \geqslant \alpha} f _ {\alpha \beta} (\mathbf {E} _ {\beta}) \neq \varnothing .
\]

Hence \(f_{\alpha}(\mathbf{E}) \neq \emptyset\) by (a), and a fortiori \(\mathbf{E} \neq \emptyset\) .

Q.E.D.

Remark. Suppose that condition (iii) in Theorem 1 is replaced by the following weaker condition:

(iii)' For each \(\alpha \in I\) and each non-empty set \(M_{\alpha} \in \mathfrak{S}_{\alpha}\) there exists \(x_{\alpha} \in M_{\alpha}\) such that \(f_{\alpha \beta}(x_{\alpha}) \in \mathfrak{S}_{\beta}\) for each \(\beta \geqslant \alpha\) .

Then the conclusion (b) of Theorem 1 remains valid; the proofs of parts (1) and (2) of Theorem 1 remain unchanged and the proof of part (3) remains

valid provided that we are careful to take \(x_{\alpha} \in \mathbf{A}_{\alpha}\) such that \(f_{\alpha \beta}(x_{\alpha})\) belongs to \(\mathfrak{S}_{\beta}\) for all \(\beta \geqslant \alpha\) . Finally, the proof of (4) shows that if

\[
\bigcap_ {\beta \geqslant \alpha} f _ {\alpha \beta} (\mathbf {E} _ {\beta}) \neq \varnothing
\]

and if we choose an \(x_{\alpha}\) in this set such that \(\overline{f_{\alpha\beta}^{-1}(x_{\alpha})} \in \mathfrak{S}_{\beta}\) whenever \(\beta \geqslant \alpha\) , there exists \(y \in E\) such that \(f_{\alpha}(y) = x_{\alpha}\) , which proves our assertion.

\minorheading{Examples}

\begin{enumerate}
\def\labelenumi{(\arabic{enumi})}
\tightlist
\item
  If the \(E_{\alpha}\) are finite sets, Theorem 1 can be applied by taking \(S_{\alpha}\) to be the set of all subsets of \(E_{\alpha}\) . * This example is generalized in General Topology to the situation in which the \(E_{\alpha}\) are compact topological spaces, the \(f_{\alpha\beta}\) continuous maps, and \(S_{\alpha}\) the set of closed subsets of \(E_{\alpha}\) (General Topology, Chapter I, §9, no. 6). *
\end{enumerate}

* (2) Let A be a ring with an identity element, and for each \(\alpha \in I\) let \(T_{\alpha}\) be an Artinian left A-module. Let \(E_{\alpha}\) be a homogeneous space for \(T_{\alpha}\) on which \(T_{\alpha}\) operates faithfully (so that \(E_{\alpha}\) is an affine space attached to \(T_{\alpha}\) ). For \(\beta \geqslant \alpha\) , suppose that \(f_{\alpha \beta}: E_{\beta} \to E_{\alpha}\) is an affine mapping. Take \(\mathfrak{S}_{\alpha}\) to be the set consisting of the empty set and the affine linear varieties in \(E_{\alpha}\) . Then condition (i) is trivially satisfied, and (ii) follows from the fact that \(T_{\alpha}\) is Artinian; for this implies that there exists a minimal element in the set of finite intersections of sets \(M \in \mathfrak{F}\) , and this minimal element must be equal to \(\bigcap_{M \in \mathfrak{F}} M\) . Finally, since \(f_{\alpha \beta}\) is affine, conditions (iii) and (iv) are trivially satisfied. *

